% ============================================================================
% 7 RESULTS
% ============================================================================
\section{Results}
\label{sec:results}

\subsection{Verified short-input baselines}

\Cref{tab:baselines} contains only rows whose saved score files match the stated
sample size. On synthetic ticket classification, a second evidence-extraction
call did not improve accuracy: zero-shot and agent two-step differ by $0.007$
on type and $0.002$ on queue, while zero-shot has the higher tag $F_1$. Both
produced valid JSON and allowed labels for every example; their quoted evidence
occurred in the input in $95.28\%$ and $94.84\%$ of cases, respectively.

\begin{table}[!t]
\caption{Verified ticket-classification results ($n=1000$). Accuracy and
macro-$F_1$ (mF$_1$) are shown for single-label fields; tags use micro-$F_1$.}
\label{tab:baselines}
\footnotesize
\setlength{\tabcolsep}{3pt}
\begin{tabular}{@{}L{1.9cm} C{0.75cm} C{0.75cm} C{0.75cm} C{0.75cm} C{0.8cm}@{}}
\toprule
\multirow{2}{*}{\textbf{Method}} & \multicolumn{2}{c}{\textbf{Type}}
 & \multicolumn{2}{c}{\textbf{Queue}} & \textbf{Tags} \\
\cmidrule(lr){2-3}\cmidrule(lr){4-5}\cmidrule(lr){6-6}
 & Acc. & mF$_1$ & Acc. & mF$_1$ & $F_1$ \\
\midrule
Zero-shot & \best{.726} & .606 & \best{.422} & \best{.366} & \best{.291} \\
Agent two-step & .719 & \best{.632} & .420 & .363 & .268 \\
\bottomrule
\end{tabular}
\end{table}

For location-only Few-NERD extraction, few-shot structured JSON reached
precision $0.7875$, recall $0.7990$ and strict $F_1=0.7932$ on 1{,}000
sentences. Few-shot free-form output reached $0.7892$, $0.7933$ and $0.7913$.
The $0.0019$ difference does not establish a format winner. On a separate
100-sentence slice, zero-shot structured and free-form reached $0.8080$ and
$0.8030$ strict $F_1$. An earlier 66-label prompt lacking an indexed token table
reached only $0.2993$ strict versus $0.6070$ lenient $F_1$ ($n=1000$), although
JSON validity was $1.0$. Since both label scope and offset instruction changed,
this diagnoses a configuration failure but does not isolate a causal prompt
effect.

\subsection{SciREX prompt selection and held-out confirmation}

On the fixed 20-paper development sample, replacing the static few-shot prompt
with the deterministic schema suggestion raised strict $F_1$ from $0.4875$ to
$0.6319$ (\Cref{tab:scirex-confirmation}). This comparison was used for prompt
selection and is not held-out evidence. The selected prompt then processed a
fresh, disjoint set of 20 official-test papers in full. Strict $F_1$ was
$0.5738$, with a document-bootstrap 95\% interval of $0.5356$--$0.6073$;
boundary-tolerant and overlap $F_1$ were $0.6736$ and $0.7030$. All 1{,}224
calls completed, with no failed document or invalid structured response.

\begin{table}[!t]
\caption{SciREX prompt development and fresh full-window confirmation. Flexible
metrics were introduced with the schema-derived prompt.}
\label{tab:scirex-confirmation}
\footnotesize
\begin{tabular}{@{}L{2.35cm} C{0.65cm} C{0.65cm} C{0.75cm} C{0.8cm} C{0.75cm}@{}}
\toprule
\textbf{Condition} & $n$ & \textbf{P} & \textbf{R} & \textbf{Strict} & \textbf{Ov.} $F_1$ \\
\midrule
Static prompt, dev prefixes & 20 & .424 & .573 & .488 & --- \\
Schema prompt, dev prefixes & 20 & .657 & .609 & \best{.632} & .742 \\
Schema prompt, held-out full & 20 & .596 & .553 & .574 & .703 \\
\bottomrule
\end{tabular}
\end{table}

On the fresh held-out sample, strict label $F_1$ was $0.5983$ for Method,
$0.5739$ for Dataset, $0.5309$ for Task and $0.4874$ for Metric. The difference
between strict and overlap scores shows that a substantial share of errors are
boundary disagreements rather than completely irrelevant suggestions. It does
not make overlap a replacement for exact scoring: exported training data still
requires the reviewer to correct those boundaries.

\subsection{100-paper full-window operational run}

The operational gate completed all 100 training papers, balanced 25 per length
bucket, and processed all 8{,}373 sentences and 9{,}901 reference mentions.
Strict micro $F_1$ was $0.5409$ (95\% paper-clustered interval
$0.5190$--$0.5622$), boundary-tolerant $F_1$ $0.6301$, and overlap $F_1$
$0.6622$. Macro-document $F_1$ was $0.5065$.

\begin{table}[!t]
\caption{Strict SciREX results by full-window length in the 100-paper
operational run. Each row contains 25 unique training papers.}
\label{tab:length}
\footnotesize
\begin{tabular}{@{}L{1.7cm} C{0.85cm} C{0.85cm} C{0.85cm} C{1.15cm}@{}}
\toprule
\textbf{Bucket} & \textbf{P} & \textbf{R} & \textbf{$F_1$} & \textbf{Overlap $F_1$} \\
\midrule
Short (1--10) & .576 & .632 & \best{.603} & .665 \\
Medium (11--50) & .537 & .567 & .552 & .665 \\
Long (51--200) & .553 & .515 & .533 & \best{.668} \\
Very long ($\geq201$) & .539 & .544 & .541 & .659 \\
\midrule
All & .544 & .538 & .541 & .662 \\
\bottomrule
\end{tabular}
\end{table}

Strict performance is lower outside the short bucket but does not decline
monotonically with window length; overlap $F_1$ stays within $0.009$ across all
four buckets. This is evidence that the runner covers long windows reliably,
not that it uses document-level context: predictions are still made one
sentence at a time.

The run made 5{,}868 logical calls in 12{,}561.81 seconds, consuming 2{,}645{,}585
input and 260{,}276 output tokens. No document failed, coverage was $100\%$, and
one structured response was invalid. The recorded cost is not reported as zero:
provider prices were not configured, so a monetary estimate would be unknown.
Per-label strict $F_1$ was $0.5919$ for Method, $0.4944$ for Task, $0.4890$ for
Metric and $0.3220$ for Dataset, making Dataset overprediction the largest
remaining label-specific problem. The prepared 1{,}000-window configuration
would require exactly 62{,}145 logical calls, but it has not been completed and
is not presented as a result.
